There is a small algebraic typo in the suggested test vector: to obtain  
\[
DCD^{-1}x=D(Cc_i),
\]
one must take  
\[
x=Dc_i,
\]
not \(x=D^{-1}c_i\). Indeed, \(D^{-1}(Dc_i)=c_i\).

\[
\textbf{Proof.}
\]

Fix an integer \(r\geq 3\), and let  
\[
D=\operatorname{diag}(d_1,\dots,d_7)
\]
with \(d_j>0\) for all \(j\). Define  
\[
A:=DCD^{-1}.
\]
Choose an index \(i\) such that  
\[
d_i=\max_{1\leq j\leq 7} d_j.
\]
Let \(c_i\) be the \(i\)th column of \(C\), and test \(A\) on the nonzero vector  
\[
x:=Dc_i.
\]
Since \(D^{-1}x=c_i\), we have
\[
Ax=DCD^{-1}(Dc_i)=D(Cc_i).
\]
By the assumed identity \(Cc_i=8e_i-\mathbf 1\), this gives
\[
Ax=D(8e_i-\mathbf 1).
\]
The vector \(8e_i-\mathbf 1\) has \(i\)th coordinate \(7\) and all other coordinates equal to \(-1\). Hence
\[
(Ax)_i=7d_i,
\qquad
(Ax)_j=-d_j \quad (j\neq i).
\]
Therefore
\[
\|Ax\|_r^r=(7d_i)^r+\sum_{j\neq i}d_j^r.
\]

Because \(C\) is a Hadamard core matrix, every entry of \(c_i\) has absolute value \(1\). Thus
\[
\|x\|_r^r
=
\|Dc_i\|_r^r
=
\sum_{j=1}^7 d_j^r |(c_i)_j|^r
=
\sum_{j=1}^7 d_j^r
=
d_i^r+\sum_{j\neq i}d_j^r.
\]

Set
\[
S:=\sum_{j\neq i}\left(\frac{d_j}{d_i}\right)^r.
\]
Since \(d_i=\max_j d_j\), we have \(0<d_j/d_i\leq 1\) for every \(j\), so
\[
0<S\leq 6.
\]
Dividing numerator and denominator by \(d_i^r\), we get
\[
\frac{\|Ax\|_r^r}{\|x\|_r^r}
=
\frac{7^r+S}{1+S}.
\]
Now compare this with \(\Gamma_r^r\). Since
\[
\Gamma_r^r=\frac{7^r+6}{7},
\]
we compute
\[
\frac{7^r+S}{1+S}-\frac{7^r+6}{7}
=
\frac{7(7^r+S)-(1+S)(7^r+6)}{7(1+S)}.
\]
Simplifying the numerator gives
\[
7(7^r+S)-(1+S)(7^r+6)
=
(7^r-1)(6-S).
\]
Hence
\[
\frac{7^r+S}{1+S}-\frac{7^r+6}{7}
=
\frac{(7^r-1)(6-S)}{7(1+S)}
\geq 0,
\]
because \(r\geq 3\) implies \(7^r-1>0\), and \(S\leq 6\). Therefore
\[
\frac{\|Ax\|_r^r}{\|x\|_r^r}
\geq
\frac{7^r+6}{7}
=
\Gamma_r^r.
\]
Since the operator norm dominates the ratio obtained from any nonzero test vector,
\[
\|A\|_{r\to r}^r
\geq
\frac{\|Ax\|_r^r}{\|x\|_r^r}
\geq
\Gamma_r^r.
\]
Taking \(r\)th roots gives
\[
\|DCD^{-1}\|_{r\to r}\geq \Gamma_r.
\]

Since \(D>0\) was arbitrary, we have shown
\[
\rho_r(C)
=
\inf_{D>0\text{ diagonal}}\|DCD^{-1}\|_{r\to r}
\geq
\Gamma_r.
\]
On the other hand, using the upper bound from the previous subproblem, namely
\[
\|C\|_{r\to r}\leq \Gamma_r,
\]
and taking the admissible diagonal matrix \(D=I_7\), we get
\[
\rho_r(C)
\leq
\|I_7CI_7^{-1}\|_{r\to r}
=
\|C\|_{r\to r}
\leq
\Gamma_r.
\]
Combining the two inequalities yields
\[
\rho_r(C)=\Gamma_r.
\]

Thus diagonal scaling cannot improve the bound.